\section{Related Work}
\label{sec:related}

\para{Sparse autoencoders for protein language models.} Interpretability of \plm{}s began with attention-based studies showing that specific heads attend to contact maps and binding sites \citep{vig2020bertology, stark2023insights}. The current wave applies \sae{}s borrowed from LLM mechanistic interpretability. \interplm \citep{simon2024interplm} trains a 10{,}240-latent ReLU \sae per \esm layer and reports that mid-layer features cover 143 of 433 \swissprot concepts at F1 $> 0.5$. Independent replication on a different model family appears in \citet{adams2025sae}. \citet{villegas2025steering} refine layer selection by intrinsic-dimension plateau and demonstrate that zinc-finger features can be steered. \citet{tsui2025lownsae} re-position \sae{}s from interpretation to prediction: a top-$k$ \sae trained on a LoRA-fine-tuned \esmlarge outperforms raw \esm in 58\% of \proteingym low-N tasks and yields top-fitness variants in 83\% of designs. \citet{protsae2025} add a semantic regularizer that trades novelty for per-feature interpretability. Unlike all of the above, our work targets the \emph{dark} tail of the same \sae inventory rather than the bright tail, and pairs three independent tests on the same feature set.

\para{Beyond feature--concept correlation.} A recent strand questions the inferential leap from feature--concept correlation to mechanistic role. \citet{haque2026antibody} compare top-$k$ and Ordered \sae{}s on autoregressive antibody \plm{}s and find that top-$k$ features identified by F1 do not reliably steer model outputs, while Ordered \sae{}s do. \citet{bricken2026actionability} argue the more general point that mechanistic-interpretability artifacts often fail to translate into actionable model edits. Concept activation vectors \citep{shamail2025cav} provide a supervised alternative: for each known motif, a linear classifier in \plm-embedding space is trained, and inner products score new windows. Our negative ablation result quantifies the \citet{haque2026antibody, bricken2026actionability} concern for ReLU \sae features in \esm specifically, and our LLM-annotation step is structurally close to the unsupervised counterpart of \citet{shamail2025cav}.

\para{LLM-assisted annotation.} \citet{parsan2025labelling} build a labelling pipeline that uses an LLM to caption \sae features and \plm neurons in ProGen2. \citet{simon2024interplm} use Claude-3.5-Sonnet to caption features and verify a glycosyltransferase cluster. We adopt the same caption strategy but constrain it to the dark tail and run a follow-up F1 re-check against the LLM-named concepts---a step neither \citet{simon2024interplm} nor \citet{parsan2025labelling} report at this granularity.

\para{Circuit-level analyses and downstream tasks.} \citet{circuit2026tracing} apply cross-layer transcoders to \esm, taking the next step beyond per-layer feature inventories. \citet{induction2026biology} identify induction-head-like circuits for repeat detection in \plm{}s, echoing the \citet{olsson2022induction} result from NLP. Variant-effect benchmarks such as \proteingym \citep{notin2023proteingym} and predecessor work \citep{notin2022tranception} provide the obvious functional anchor for \sae feature claims. Our \proteingym experiment is in this lineage; we use the per-residue mean $|\dms|$ score across three assays.

\para{Position of our work.} We sit at the intersection of three lines: (1) the \plm-\sae feature-inventory work \citep{simon2024interplm, adams2025sae, villegas2025steering}, (2) the actionability critique \citep{haque2026antibody, bricken2026actionability}, and (3) LLM-assisted annotation \citep{parsan2025labelling}. To our knowledge we are the first to triangulate all three on the same feature set with a quantitative null model, and the first to test whether ``low-F1 \emph{dark}'' is the right operationalization of ``unknown biology'' for the SAE feature inventory of a public \plm.
